%% toc: P2 19 19.8
\section{Factoring under translation and superposition into either side}
\label{sec:p02-w02c-p02-021-two-steps}

\Cref{sec:p02-w02b-p02-026-equality} connects clause sets with equality to
the equality-free resolution calculus through the translation
\((\cdot)^{\mathsf{eq}}\) of Convention~\ref{def:p02-w02b-p02-026-reading}, and
proves the soundness of the ordered-superposition rules.  Two steps of that
development are stated and proved here as lemmas of their own: the
correspondence of single Factoring steps under translation, in both
directions, and the soundness of a Superposition step whose replaced
occurrence lies in the second argument of the target literal.

\emph{Purpose.}  The first lemma completes the converse direction of
Lemma~\ref{thm:p02-w02b-p02-026-translation}(3), on which the
equivalence of (2) and (4) in
Theorem~\ref{thm:p02-w02b-p02-026-eq-resolution} and the proof of
Theorem~\ref{thm:p02-w02b-p02-026-bare} rest.  The second completes the case
analysis in the proof of
Lemma~\ref{thm:p02-w02b-p02-026-superposition-sound}.
\emph{Inputs.}  \(\Lang\)-clauses in the sense of
\cref{sec:fol-cardinality-normalforms-04}, whose literals may be equations;
the fresh relation symbol \(\mathsf{eq}\notin\Lang\); for the second lemma, the
premises of a Superposition inference of the rule sheet of
Definition~\ref{def:p12-ordered-superposition}, read as in
\cref{sec:p02-w02b-p02-026-superposition}.
\emph{Conclusions.}  Lemma~\ref{thm:p02-w02c-p02-021-factoring} and
Lemma~\ref{thm:p02-w02c-p02-021-superposition-rhs}.  All proofs are in the
annex.

\subsection{Factoring steps under translation}
\label{sec:p02-w02c-p02-021-factoring}

Recall the Factoring rule of \cref{sec:fol-cardinality-normalforms-04}: from
a clause \(E\) containing two distinct literals of the same polarity whose
atoms \(A\) and \(B\) have a most general unifier \(\theta\), it infers
\(E\theta\).  Say that \(C'\) \emph{follows from \(C\) by Factoring on
\(L_1,L_2\) with \(\theta\)} when \(L_1,L_2\) are such literals of \(C\),
\(\theta\) is such a unifier of their atoms, and \(C'=C\theta\).  Nieuwenhuis
and Rubio display the same rule beside binary resolution, with the two
merged literals written \(A\) and \(B\) and the conclusion
\((C\lor A)\sigma\) for \(\sigma=\mathit{mgu}(A,B)\)
\cite[\S1.1, p.~3]{NieuwenhuisRubio2001Typescript}.

\begin{hblemma}{thm:p02-w02c-p02-021-factoring}{proved=annex:p02-w02c-p02-021-factoring}{Factoring commutes with translation}
Let \(C\) and \(C'\) be \(\Lang\)-clauses and \(\theta\) a term substitution.
\begin{enumerate}
  \item If \(C'\) follows from \(C\) by Factoring on \(L_1,L_2\) with
        \(\theta\), then \(C'^{\mathsf{eq}}\) follows from \(C^{\mathsf{eq}}\)
        by Factoring on \(L_1^{\mathsf{eq}},L_2^{\mathsf{eq}}\) with
        \(\theta\).
  \item If \(C'^{\mathsf{eq}}\) follows from \(C^{\mathsf{eq}}\) by Factoring
        on \(M_1,M_2\) with \(\theta\), then there are distinct literals
        \(L_1,L_2\in C\) of the same polarity with
        \(L_1^{\mathsf{eq}}=M_1\) and \(L_2^{\mathsf{eq}}=M_2\), \(\theta\) is
        a most general unifier of their atoms, and \(C'=C\theta\); that is,
        \(C'\) follows from \(C\) by Factoring on \(L_1,L_2\) with \(\theta\).
\end{enumerate}
Consequently Lemma~\ref{thm:p02-w02b-p02-026-translation}(3) holds in both
directions for derivations that use Factoring.
\end{hblemma}

The proof of item~2 uses that \(\mathsf{eq}\) does not belong to \(\Lang\).
Without that hypothesis the translation can identify distinct literals, and
the converse fails.

\begin{hbwarning}{def:p02-w02c-p02-021-fresh}{The fresh symbol is needed for the converse}
Suppose, contrary to Convention~\ref{def:p02-w02b-p02-026-reading}, item~3,
that \(\Lang\) already contained a binary relation symbol \(\mathsf{eq}\), and
let the translation still send \(r=t\) to \(r\mathrel{\mathsf{eq}}t\).  Take a
constant \(a\) and variables \(x\neq y\), and
\(C=\{x=a,\ y\mathrel{\mathsf{eq}}a\}\).  Then
\(C^{\mathsf{eq}}=\{x\mathrel{\mathsf{eq}}a,\ y\mathrel{\mathsf{eq}}a\}\), and
Factoring on its two positive literals with the most general unifier
\(\{x\mapsto y\}\) gives \(\{y\mathrel{\mathsf{eq}}a\}\).  Every
\(\Lang\)-clause \(C'\) with \(C'^{\mathsf{eq}}=\{y\mathrel{\mathsf{eq}}a\}\)
is \(\{y=a\}\), \(\{y\mathrel{\mathsf{eq}}a\}\) or
\(\{y=a,\ y\mathrel{\mathsf{eq}}a\}\), yet the atoms \(x=a\) and
\(y\mathrel{\mathsf{eq}}a\) of \(C\) have different atomic symbols and do not
unify, so no Factoring step from \(C\) yields any clause.  Item~2 of
Lemma~\ref{thm:p02-w02c-p02-021-factoring} fails, and with it the reading of
derivations from \(T^{\mathsf{eq}}\) as derivations from \(T\).
\end{hbwarning}

\begin{hbexample}{def:p02-w02c-p02-021-factoring-example}{Factoring an equation and its translation}
Let \(a\) be a constant and \(C=\{x=a,\ a=y\}\) with \(x\neq y\).  The
procedure of Lemma~\ref{thm:p02-w02-p02-028-unification} applied to the atoms
\(x=a\) and \(a=y\) decomposes them into the argument pairs \((x,a)\) and
\((a,y)\) and returns \(\theta=\{x\mapsto a,\ y\mapsto a\}\).  Factoring
gives \(C\theta=\{a=a\}\), the two literals having become one member.  On the
translated side, \(C^{\mathsf{eq}}=\{x\mathrel{\mathsf{eq}}a,\
a\mathrel{\mathsf{eq}}y\}\), the same pairs arise, and Factoring with the same
\(\theta\) gives \(\{a\mathrel{\mathsf{eq}}a\}=(C\theta)^{\mathsf{eq}}\), as
item~1 states.  Conversely, starting from the translated step, item~2
recovers \(L_1=(x=a)\) and \(L_2=(a=y)\) as the only literals of \(C\) whose
translations are the two factored literals, and returns the clause
\(\{a=a\}\).
\end{hbexample}

\subsection{Superposition into the second argument of the target literal}
\label{sec:p02-w02c-p02-021-superposition}

As in \cref{sec:p02-w02b-p02-026-superposition}, an equation \(s\approx t\)
of the rule sheet is the atom \(s=t\) and \(s\not\approx t\) is
\(\neg(s=t)\); the two arguments of an atom are ordered.  The sheet's
Superposition rule displays its distinguished occurrence \(l'\) in the first
argument of the target literal, \(s[l']\mathrel{\diamond}t\).  The case in
which the occurrence lies in the second argument, \(t\mathrel{\diamond}s[l']\),
is the following; it also allows the equation of the first premise to be
written with its sides in either order.  In the chapter of Nieuwenhuis and
Rubio this case does not arise separately, because there an equation is the
multiset \(\{s,t\}\), so that \(s\simeq t\) and \(t\simeq s\) denote the same
equation \cite[\S2.1, p.~12]{NieuwenhuisRubio2001Typescript}, and a rewrite
step may use an equation \(l\simeq r\) in either direction
\cite[\S1.2, p.~5]{NieuwenhuisRubio2001Typescript}.  With ordered
arguments, the lemma below is what that identification amounts to.

\begin{hblemma}{thm:p02-w02c-p02-021-superposition-rhs}{proved=annex:p02-w02c-p02-021-superposition-rhs}{Superposition into the second argument is sound}
Let \(C\) and \(D\) be \(\Lang\)-clauses, \(l,r,s,t\) terms, \(p\) a position
of \(s\) with \(s|_p=l'\), \(\diamond\) either \(=\) or its negation, and
\(\sigma\) a substitution with \(l\sigma=l'\sigma\).  Let \(K\) be \(l=r\) or
\(r=l\).  Every \(\Lang\)-structure that satisfies the premises
\(C\cup\{K\}\) and \(D\cup\{t\mathrel{\diamond}s\}\) satisfies the conclusion
\[
 \bigl(C\cup D\cup\{t\mathrel{\diamond}s[r]_p\}\bigr)\sigma .
\]
The same holds with \(s\mathrel{\diamond}t\) and \(s[r]_p\mathrel{\diamond}t\)
in place of \(t\mathrel{\diamond}s\) and \(t\mathrel{\diamond}s[r]_p\).  The
proof uses only that \(\sigma\) unifies \(l\) and \(l'\); the ordering,
eligibility and selection conditions of the sheet are not used.
\end{hblemma}

With this lemma every Superposition inference of the sheet, wherever its
distinguished occurrence lies and however its equations are oriented, is
sound, and Lemma~\ref{thm:p02-w02b-p02-026-superposition-sound} holds as
stated.  Soundness is the only property proved here; refutational
completeness of the sheet is the cited
Theorem~\ref{thm:p12-superposition-completeness}.

\begin{hbexample}{def:p02-w02c-p02-021-superposition-example}{A replacement in the second argument}
Let \(c\) be a constant and \(f,h\) unary function symbols.  Take the premises
\(\{f(x)=x\}\) and \(\{\neg(c=h(f(c)))\}\), so \(C=D=\varnothing\), \(l=f(x)\),
\(r=x\), \(t=c\), \(s=h(f(c))\), \(\diamond\) is the negation of \(=\), and the
occurrence \(l'=f(c)\) sits at position \(1\) of \(s\), in the second argument
of the target literal.  The substitution \(\sigma=\{x\mapsto c\}\) gives
\(l\sigma=f(c)=l'\sigma\), and the conclusion is
\(\{\neg(c=h(c))\}\).  In any structure satisfying the premises,
\(f(c)\) and \(c\) denote the same element, so \(h(f(c))\) and \(h(c)\) do,
and \(c\) differs from both.
\end{hbexample}

\emph{Relationships.}  Lemma~\ref{thm:p02-w02c-p02-021-factoring} and the
Resolution case written in the proof of
Lemma~\ref{thm:p02-w02b-p02-026-translation} together give item~(3) of that
lemma in both directions; Theorem~\ref{thm:p02-w02b-p02-026-eq-resolution},
the equality-complete connection of \cref{sec:p02-w02b-p02-026-method}, and
Theorem~\ref{thm:p02-w02b-p02-026-bare} use that item in both directions.
Lemma~\ref{thm:p02-w02c-p02-021-superposition-rhs} supplies the remaining case
of Lemma~\ref{thm:p02-w02b-p02-026-superposition-sound}.

\emph{Sources.}  The Resolution and Factoring rules are those of
\cref{sec:fol-cardinality-normalforms-04}; the Superposition rule is that of
the rule sheet of Definition~\ref{def:p12-ordered-superposition}; positions
and the facts about them are those set out at the start of the proof of
Lemma~\ref{thm:p02-w02b-p02-026-superposition-sound}.  The displayed
factoring rule is \cite[\S1.1, p.~3]{NieuwenhuisRubio2001Typescript}; the
reading of an equation as an unordered pair is
\cite[\S2.1, p.~12]{NieuwenhuisRubio2001Typescript}, with use in either
direction at \cite[\S1.2, p.~5]{NieuwenhuisRubio2001Typescript}; page
locators refer to the typescript version of that chapter.  The proofs are
written for this handbook.
